\documentclass[12pt,a4paper]{article}
\usepackage[utf8]{inputenc}
\usepackage[T2A]{fontenc}
\usepackage[russian]{babel}
\usepackage{amsmath,amssymb,amsthm}
\usepackage[margin=2.5cm]{geometry}
\usepackage[unicode,hidelinks]{hyperref}
\setlength{\emergencystretch}{3em}
\newtheorem{theorem}{Теорема}
\newtheorem{lemma}{Лемма}
\newtheorem{corollary}{Следствие}
\theoremstyle{definition}
\newtheorem{remark}{Замечание}
\renewcommand{\qedsymbol}{}
\newcommand{\step}[1]{\smallskip\noindent\textbf{Шаг #1.}}
\newcommand{\R}{\mathbb R}
\renewcommand{\C}{\mathbb C}
\newcommand{\norm}[1]{\lVert#1\rVert}
\newcommand{\abs}[1]{\lvert#1\rvert}
\newcommand{\inner}[2]{\langle#1,#2\rangle}
\title{О базисности корневых функций оператора второго порядка\\
с отражением аргумента в старшей производной}
\author{Л.\,И.~Дядченко}
\date{}
\begin{document}
\maketitle
\begin{abstract}
Рассматривается оператор $L_ay=-y''(x)-a(x)y''(-x)$ в пространстве $L_2(-1,1)$ с
краевыми условиями Дирихле. Здесь $a$ --- вещественная ограниченная функция, для которой
$\abs{a(x)a(-x)}\le\kappa<1$. Доказаны компактность обратного оператора, полнота системы
корневых функций для непрерывного $a$, базисность Рисса системы собственных функций для
чётного $a$ и базисность Рисса со скобками системы корневых функций для знакопостоянного
$a$ из $W^2_\infty$. Если чётная часть $a$ по модулю меньше единицы, корневые функции
образуют базис для метода суммирования Абеля--Лидского любого порядка из интервала,
который определяется отношением нечётной части к чётной. Построены примеры, в которых оценка произведения выполнена, но
базиса Рисса из собственных функций нет; среди них $a(x)=x/2$. Результаты дополняют
известные теоремы о постоянном коэффициенте и опираются на теорию подчинённых
возмущений.
\end{abstract}

\section{Введение}
В пространстве $H=L_2(-1,1)$ комплекснозначных функций изучается оператор
\begin{equation}\label{eq:op}
 L_ay=-y''(x)-a(x)\,y''(-x),\qquad
 D(L_a)=\{y\in W_2^2(-1,1):\ y(-1)=y(1)=0\}.
\end{equation}
Здесь $y''(-x)$ --- значение второй производной в точке $-x$, а $a$ --- вещественная
функция из $L_\infty(-1,1)$, для которой при некотором $\kappa\in[0,1)$ почти всюду
\begin{equation}\label{eq:bound}
 \abs{a(x)\,a(-x)}\le\kappa .
\end{equation}

Всюду далее $(Rf)(x)=f(-x)$ --- оператор отражения, $A_0$ --- оператор $-d^2/dx^2$ с
областью $D(A_0)=D(L_a)$ и $G=A_0^{-1}$. Функцию и оператор умножения на неё обозначаем
одной буквой. Положим $T_a=I+aR$; тогда $L_a=T_aA_0$. Напомним, что оператор $G$
компактен, самосопряжён и положителен, а оператор $R$ унитарен, $R^2=I$, и коммутирует с
$A_0$ и $G$. Собственные и присоединённые функции называем корневыми.

Спектральным задачам для уравнений с отражением аргумента посвящено много работ, из
которых отметим \cite{Pol,KS}. Для постоянного коэффициента $a\equiv c$
в~\cite[теорема~4.1]{KS} доказано, что собственные функции задачи Дирихле образуют
ортонормированный базис. В~\cite[теорема~1.1]{Pol} для операторов чётного порядка выше
второго с отражением в старшей производной методом подобных операторов доказана
базисность Бари. Однако в этих работах коэффициент при отражённой старшей производной
постоянен. При переменном $a$ отражение входит в старшую часть оператора, и $L_a$ не
является возмущением оператора $A_0$ младшими членами.

Цель работы --- выяснить, при каких условиях на коэффициент $a$ система корневых функций
оператора $L_a$ полна, при каких образует базис Рисса и при каких --- базис Рисса со
скобками. Основной результат получен в теореме~\ref{th:sign}: если
$a\in W^2_\infty(-1,1)$ и $\abs a\ge m>0$, то система корневых функций полна и образует
базис Рисса со скобками. Для чётного $a$ собственные функции образуют обычный базис Рисса
(теорема~\ref{th:even}). Для непрерывного $a$ система корневых функций полна
(теорема~\ref{th:complete}). Если $\norm{a_+}_\infty<1$, где $a_+$ --- чётная часть
$a$, то корневые функции образуют базис для метода суммирования Абеля--Лидского порядка
$\beta$ при каждом $\beta\in(1/2,\tau^{-1})$; здесь $\pi\tau/2$ --- угол сектора, в
котором лежат собственные значения (теорема~\ref{th:abel}). Знак $a$ при этом не
важен. В теоремах~\ref{th:onesided}, \ref{th:ramp}
и~\ref{th:half} показано, что одного условия~\eqref{eq:bound} для базисности
недостаточно: базиса Рисса из собственных функций нет уже при $a(x)=x/2$.

Метод исследования таков. Оператор $T_a$ ограниченно обратим, поэтому $L_a$ имеет
компактный обратный (\S2). При чётном $a$ оператор $L_a$ подобен самосопряжённому (\S3).
В условиях теоремы~\ref{th:sign} умножение на подходящий вес приводит $L_a$ к оператору с
чётным коэффициентом, возмущённому оператором первого порядка, и применима теорема
Маркуса--Мацаева о подчинённых возмущениях \cite[теорема~6.1]{Wyss},
\cite[теорема~6.6]{Shk} (\S4). Для коэффициента, близкого к чётному, эта теорема
заменяется прямым построением (\S5). Полнота в теореме~\ref{th:complete} выводится из
оценки резольвенты на отрицательной полуоси и теоремы о полноте
\cite[теорема~4.7]{Shk}; оценку даёт умножение на вес $e^{\theta/2}$ с нечётной функцией
$\theta$, после которого чётная часть коэффициента по модулю меньше единицы (\S6). Если
чётная часть мала с самого начала, резольвента убывает как $\abs\lambda^{-1}$ на всех
лучах вне сектора собственных значений, и суммируемость по Абелю--Лидскому в
теореме~\ref{th:abel} следует из теоремы Лидского--Мацаева \cite[теорема~5.1]{Shk}. Для
коэффициентов, меняющих знак, доказаны только полнота и суммируемость; вопрос о
базисности со скобками требует иного подхода и здесь не рассматривается.

Напомним определения. Последовательность замкнутых подпространств $\{W_k\}$ пространства
$H$ назовём \emph{базисом Рисса подпространств}, если её линейная оболочка плотна в $H$ и
найдутся постоянные $0<\alpha_1\le\alpha_2$, для которых
\[
 \alpha_1\sum_k\norm{v_k}^2\le\Bigl\lVert\sum_kv_k\Bigr\rVert^2\le\alpha_2\sum_k\norm{v_k}^2
\]
при любом конечном наборе векторов $v_k\in W_k$. Будем говорить, что корневые функции
оператора $L$ образуют \emph{базис Рисса со скобками}, если существует базис Рисса
подпространств $\{W_k\}$, состоящий из конечномерных инвариантных относительно $L$
подпространств области $D(L)$, каждое из которых порождено корневыми функциями $L$. Если
все $W_k$ одномерны, то нормированные корневые функции образуют обычный базис Рисса.

\section{Обратимость оператора $T_a$}
\begin{lemma}\label{lem:inverse}
Пусть выполнено условие~\eqref{eq:bound}. Тогда оператор $T_a$ ограниченно обратим в $H$ и
\[
 T_a^{-1}=d\,(I-aR),\qquad d(x)=\frac1{1-a(x)a(-x)}.
\]
\end{lemma}
\begin{proof}
Заметим, что $(aR)(aR)f(x)=a(x)\,a(-x)\,f(x)$. Поэтому
\[
 (I-aR)(I+aR)=(I+aR)(I-aR)=1-a(x)a(-x)\ge1-\kappa .
\]
Функция $d$ чётна и ограничена числом $1/(1-\kappa)$, а потому умножение на $d$
коммутирует с $R$ и с умножением на $a$. Лемма доказана.
\end{proof}

\begin{theorem}\label{th:class}
Пусть выполнено условие~\eqref{eq:bound}. Тогда оператор $K_a=GT_a^{-1}$ компактен,
инъективен и $K_a=L_a^{-1}$. Более того, собственные подпространства оператора $L_a$
конечномерны, и для каждого $\lambda\in\C$ либо $\lambda$ --- собственное значение $L_a$,
либо уравнение $(L_a-\lambda)f=g$ однозначно разрешимо в $D(L_a)$ при любом $g\in H$.
\end{theorem}
\begin{proof}
Из леммы~\ref{lem:inverse} получаем $L_aK_a=T_aA_0GT_a^{-1}=I$ и $K_aL_a=I$ на
$D(L_a)$. Оператор $K_a$ компактен вместе с $G$. При $\lambda\ne0$ равенство
$L_af=\lambda f$ равносильно равенству $K_af=\lambda^{-1}f$, и остальные утверждения
следуют из альтернативы Фредгольма для $K_a$. Теорема доказана.
\end{proof}

\begin{remark}\label{rem:matrix}
В координатах $(y(x),y(-x))$ на интервале $(0,1)$ оператор $T_a$ есть умножение на
матрицу
\[
 M(x)=\begin{pmatrix}1&a(x)\\a(-x)&1\end{pmatrix},\qquad\det M(x)=1-a(x)a(-x),
\]
и лемма~\ref{lem:inverse} выражает её обратимость. Отметим, что при $\kappa=1$
теорема~\ref{th:class} неверна: для $a\equiv1$ все функции $\sin(m\pi x)$ лежат в ядре
$L_a$.
\end{remark}

\begin{lemma}\label{lem:sym}
Оператор $T_a$ симметричен в $H$ в том и только том случае, когда $a(-x)=a(x)$ почти
всюду.
\end{lemma}
\begin{proof}
Имеем $(aR)^*=Ra$. Поэтому симметричность равносильна равенству
$a(x)g(-x)=a(-x)g(-x)$ для всех $g\in H$, а оно при $g\equiv1$ даёт чётность $a$.
Лемма доказана.
\end{proof}
Существо различия между чётным и общим случаем состоит в следующем: при чётном $a$
оператор $T_a$ положителен, и $L_a=T_aA_0$ подобен самосопряжённому оператору
$T_a^{1/2}A_0T_a^{1/2}$; при нечётной части, отличной от нуля, такого подобия лемма не
даёт.

\section{Чётный коэффициент}
\begin{theorem}\label{th:even}
Пусть выполнено условие~\eqref{eq:bound} и $a(-x)=a(x)$ почти всюду. Тогда собственные
значения оператора $L_a$ положительны, и существует базис Рисса пространства $H$,
составленный из его собственных функций. Более того, присоединённых функций нет: если
$f,g\in D(L_a)$, $(L_a-\lambda)f=0$ и $(L_a-\lambda)g=f$, то $f=0$.
\end{theorem}
\begin{proof}
Доказательство разобьём на три шага.

\step{1} Из~\eqref{eq:bound} и чётности следует $\abs a\le\rho=\sqrt\kappa<1$. Обозначим
через $H_s$, $s=\pm1$, подпространство функций с $Rf=sf$ и положим $p_s=1+sa$; тогда
$1-\rho\le p_s\le1+\rho$. Если $f\in D(L_a)\cap H_s$, то $f''\in H_s$ и
\[
 L_af=-f''-a\,Rf''=-(1+sa)f''=p_sA_0f .
\]

\step{2} Положим $w=p_s^{-1/2}$. Функция $w$ чётна, поэтому умножение на неё сохраняет
$H_s$, и оператор $wGw$ в $H_s$ компактен, самосопряжён и положителен. Пусть $\{u_i\}$
--- ортонормированный базис $H_s$ из его собственных векторов, $wGwu_i=\mu_iu_i$,
$\mu_i>0$. Функция $y_i=G(wu_i)$ лежит в $D(L_a)\cap H_s$ и равна
$\mu_ip_s^{1/2}u_i$, откуда
\[
 L_ay_i=p_sA_0y_i=p_swu_i=p_s^{1/2}u_i=\mu_i^{-1}y_i .
\]
Система $\{p_s^{1/2}u_i\}$ есть образ ортонормированного базиса при ограниченно обратимом
операторе умножения, то есть базис Рисса в $H_s$. Подпространства $H_+$ и $H_-$
ортогональны и в сумме дают $H$, поэтому объединение двух систем --- базис Рисса в $H$.

\step{3} Проекторы $(I\pm R)/2$ сохраняют $D(L_a)$ и коммутируют с $L_a$, поэтому
оставшиеся утверждения достаточно доказать в $H_s$. В эквивалентном скалярном
произведении $[f,g]=\int_{-1}^1f\bar g\,p_s^{-1}dx$ для $f,g\in D(L_a)\cap H_s$ имеем
\[
 [L_af,g]=\int_{-1}^1f'\,\overline{g'}\,dx=[f,L_ag] .
\]
В частности, $[L_af,f]>0$ при $f\ne0$, и собственные значения положительны. Если
$(L_a-\lambda)f=0$, $f\ne0$ и $(L_a-\lambda)g=f$, то $\lambda$ вещественно и
$[f,f]=[(L_a-\lambda)g,f]=[g,(L_a-\lambda)f]=0$, что невозможно. Теорема доказана.
\end{proof}

\begin{remark}\label{rem:const}
При $a\equiv c\in(-1,1)$ получаем ортогональный базис из функций
$\cos((n+\tfrac12)\pi x)$, $n\ge0$, с собственными значениями
$(1+c)(n+\tfrac12)^2\pi^2$ и функций $\sin(m\pi x)$, $m\ge1$, с собственными значениями
$(1-c)m^2\pi^2$, что согласуется с~\cite[теорема~4.1]{KS}. Отметим, что гладкость $a$ в
теореме~\ref{th:even} не используется.
\end{remark}

\section{Коэффициент постоянного знака}
Сформулируем основной результат работы.
\begin{theorem}\label{th:sign}
Пусть выполнены следующие условия:
\begin{enumerate}
\item[(i)] функция $a$ вещественна и принадлежит $W^2_\infty(-1,1)$;
\item[(ii)] $0<m\le\abs{a(x)}\le M$ почти всюду (а потому $a$ имеет постоянный знак);
\item[(iii)] выполнено неравенство~\eqref{eq:bound}.
\end{enumerate}
Тогда существует последовательность конечномерных подпространств $W_k\subset D(L_a)$,
инвариантных относительно $L_a$ и образующих базис Рисса подпространств в $H$. Более
того, каждое $W_k$ порождено корневыми функциями оператора $L_a$; в частности, система
корневых функций полна в $H$ и образует базис Рисса со скобками.
\end{theorem}
Чётность $a$ не предполагается; условиям теоремы удовлетворяет, например,
$a(x)=\tfrac12+\tfrac x4$ с $\kappa=\tfrac14$.

Идея доказательства состоит в следующем. Умножение на вес $r=\sqrt{\abs a}$ переводит
$L_a$ в оператор $L_b+S$ с чётным коэффициентом $b$ и оператором первого порядка $S$
(лемма~\ref{lem:gauge}). По теореме~\ref{th:even} оператор $L_b$ подобен диагональному
оператору $\Lambda$; возмущение подчинено ему с показателем $1/2$
(лемма~\ref{lem:sub}), а собственные значения $\Lambda$ растут квадратично
(лемма~\ref{lem:count}). В этих условиях применима теорема о подчинённых возмущениях
(теорема~\ref{th:mm}). Отметим, что сразу применить её к $L_a$ затруднительно, так как
отражение входит в старшую часть оператора. Теперь реализуем намеченный план.

\begin{lemma}\label{lem:gauge}
Пусть $r$ и $1/r$ принадлежат $W^2_\infty(-1,1)$, $a,b\in L_\infty(-1,1)$ и
$b(x)r(x)=a(x)r(-x)$ почти всюду. Тогда $T_ar=rT_b$ и для $f\in D(A_0)$
\begin{equation}\label{eq:gauge}
 L_a(rf)=r\bigl(L_bf+T_bB_0f\bigr),\qquad
 B_0f=-2\,\frac{r'}{r}\,f'-\frac{r''}{r}\,f .
\end{equation}
\end{lemma}
\begin{proof}
Имеем $(I+aR)(rf)(x)=r(x)f(x)+a(x)r(-x)f(-x)=r(x)\bigl(f(x)+b(x)f(-x)\bigr)$. Далее,
$-(rf)''=r\,(A_0f+B_0f)$, и остаётся применить $T_a$. Умножение на $r$ и на $1/r$
сохраняет $D(A_0)$. Лемма доказана.
\end{proof}
В условиях теоремы~\ref{th:sign} функция $a$ непрерывна и отделена от нуля; обозначим
через $\sigma$ её знак и положим $r=\sqrt{\abs a}$. Тогда
\[
 b(x)=\sigma\sqrt{a(x)\,a(-x)}
\]
--- чётная функция, $\abs{b(x)b(-x)}=\abs{a(x)a(-x)}\le\kappa$, и в силу~\eqref{eq:gauge}
оператор $L_a$ подобен оператору $L_b+S$, где $S=T_bB_0$.

Обозначим $P_\pm=(I\pm R)/2$, $p_\pm=1\pm b$ и положим
\[
 V=p_+^{1/2}P_++p_-^{1/2}P_-,\qquad W=V^{-1}=p_+^{-1/2}P_++p_-^{-1/2}P_-,\qquad
 K=WGW,\qquad\Lambda=K^{-1}.
\]
Операторы $V$ и $W$ ограничены и самосопряжены, причём
$\norm W\le\nu:=(1-\sqrt\kappa)^{-1/2}$. Согласно шагам~1 и~2 доказательства
теоремы~\ref{th:even} оператор $K$ компактен, самосопряжён и положителен, и
$L_b=V\Lambda V^{-1}$. Пусть $\{e_i\}$ --- ортонормированный базис из собственных
векторов $\Lambda$ и $\Lambda e_i=\lambda_ie_i$.

\begin{lemma}\label{lem:sub}
Оператор $\tilde S=V^{-1}SV$ определён на $D(\Lambda)$ и
\begin{equation}\label{eq:sub}
 \norm{\tilde Su}\le b_1\inner{\Lambda u}{u}^{1/2}+b_0\norm u,\qquad
 b_1=2\nu(1+\sqrt\kappa)\,\Bigl\lVert\frac{r'}r\Bigr\rVert_\infty ,
\end{equation}
с некоторой постоянной $b_0$.
\end{lemma}
\begin{proof}
Пусть $u\in D(\Lambda)$. Тогда функция $f=Vu$ лежит в $D(A_0)$ и
$\inner{\Lambda u}{u}=\inner{VA_0f}{Wf}=\inner{A_0f}{f}=\norm{f'}^2$. Следовательно,
\[
 \norm{\tilde Su}\le\norm W\,\norm{T_b}\Bigl(2\Bigl\lVert\frac{r'}r\Bigr\rVert_\infty
 \norm{f'}+\Bigl\lVert\frac{r''}r\Bigr\rVert_\infty\norm f\Bigr),
\]
и остаётся учесть оценки $\norm{T_b}\le1+\sqrt\kappa$ и $\norm f\le\norm V\norm u$.
Лемма доказана.
\end{proof}

\begin{lemma}\label{lem:count}
При каждом $t>0$ число собственных значений $\lambda_i\le t^2$ с учётом кратности не
превосходит $\dfrac{2\nu}{\pi}\,t+2$.
\end{lemma}
\begin{proof}
Собственные значения оператора $G$ равны $(2/(k\pi))^2$, $k\ge1$. Обозначим через $F_N$
часть спектрального разложения $G$ с номерами $k\le2N$ и положим $E_N=G-F_N$; тогда
$\operatorname{rank}F_N\le2N$ и $\norm{E_N}<(N\pi)^{-2}$. Выберем наименьшее натуральное
$N$, для которого $\nu^2(N\pi)^{-2}\le t^{-2}$. Пусть нашлось более $2N$ ортонормированных
собственных векторов оператора $K$ с собственными значениями
$\lambda_i^{-1}\ge t^{-2}$. Тогда в их линейной оболочке есть ненулевой вектор $u$ с
$F_NWu=0$, и
\[
 t^{-2}\norm u^2\le\inner{Ku}{u}=\inner{E_NWu}{Wu}<\nu^2(N\pi)^{-2}\norm u^2 ,
\]
что противоречит выбору $N$. Значит, таких векторов не больше $2N\le2\nu t/\pi+2$.
Лемма доказана.
\end{proof}

Напомним теорему о подчинённых возмущениях в нужной нам форме.
\begin{theorem}\label{th:mm}
Пусть оператор $\Lambda$ диагонален в ортонормированном базисе $\{e_i\}$ гильбертова
пространства, $\Lambda e_i=\lambda_ie_i$, $\lambda_i>0$, и выполнены следующие условия:
\begin{enumerate}
\item[(i)] число значений $\lambda_i\le t^2$ не превосходит $c_2t+C_0$ при всех $t\ge0$;
\item[(ii)] оператор $S$ определён на $D(\Lambda)$ и
$\norm{Su}\le b_1\inner{\Lambda u}{u}^{1/2}+b_0\norm u$.
\end{enumerate}
Тогда существуют конечномерные подпространства $V_k\subset D(\Lambda)$, инвариантные
относительно $\Lambda+S$ и образующие базис Рисса подпространств.
\end{theorem}
Это утверждение --- частный случай теоремы Маркуса--Мацаева в
форме~\cite[теорема~6.1]{Wyss}; см. также \cite[теорема~6.6]{Shk}. Действительно,
оператор $\Lambda$ самосопряжён, положителен и имеет компактную резольвенту. Из условия
(ii) и неравенств $\inner{\Lambda u}{u}\le\norm{\Lambda u}\norm u$,
$\norm u\le\lambda_{\min}^{-1}\norm{\Lambda u}$ следует
$\norm{Su}\le b\,\norm{\Lambda u}^{1/2}\norm u^{1/2}$, то есть $S$ подчинён $\Lambda$ с
показателем $p=1/2$. Наконец, в силу (i) число собственных значений $\lambda_i\le r$ есть
$O(r^{1-p})$. В обозначениях~\cite{Shk} порядок оператора равен $\alpha=2$, и условие
$\alpha^{-1}\le1-p$ выполнено как равенство. Для удобства читателя приведём схему
доказательства.

\begin{proof}[Схема доказательства]
\step{1} Заменим $\Lambda$ на $\Lambda_c=\Lambda+c$ и положим
$\omega_i=\sqrt{\lambda_i+c}$. Инвариантные подпространства от сдвига не зависят, а
оператор $\hat S=S\Lambda_c^{-1/2}$ ограничен и
$\beta:=\norm{\hat S}\le b_1+b_0/\sqrt c$.

\step{2} Для $T=\Lambda_c+S$ в переменной $\omega=\sqrt z$ имеем
\[
 T-\omega^2=\bigl(I+N(\omega)\bigr)(\Lambda_c-\omega^2),\qquad
 N(\omega)=\hat S\,\Lambda_c^{1/2}(\Lambda_c-\omega^2)^{-1},
\]
и при $\operatorname{Re}\omega\ge0$
\[
 \norm{N(\omega)}\le\beta\sup_i\frac{\omega_i}{\abs{\omega_i^2-\omega^2}}\le
 \frac{\beta}{\operatorname{dist}(\omega,\{\omega_i\})}.
\]
Если это расстояние не меньше $2\beta$, то оператор $I+N(\omega)$ обратим. В частности,
спектр $T$ лежит в полосе $\abs{\operatorname{Im}\omega}<2\beta$.

\step{3} Строим вертикальные отрезки $\operatorname{Re}\omega=\xi_k$,
$\abs{\operatorname{Im}\omega}\le\alpha$, на которых норма $(I+N(\omega))^{-1}$ ограничена
равномерно по $k$. Из условия (i) по принципу Дирихле следует, что сколь угодно далеко
есть окна фиксированной ширины, содержащие не более $M_0$ частот $\omega_i$. Внутри окна
частоты могут сближаться, и оценка шага~2 не работает. Отщепим их:
$I+N(\omega)=A(\omega)+F(\omega)$, где $A(\omega)$ отвечает частотам вне окна,
$\norm{A(\omega)^{-1}}\le2$, а ранг $F(\omega)$ не превосходит $M_0$. По формуле Вудбери
обратимость $I+N(\omega)$ равносильна обратимости матрицы порядка не выше $M_0$. Её
определитель после нормировки --- голоморфная функция $\omega$, ограниченная в окне
постоянной, не зависящей от окна, и отделённая от нуля в точке высоко над осью, где
применим шаг~2. Число её нулей ограничено; после деления на нули теорема
Бореля--Каратеодори для логарифма частного даёт оценку снизу на вертикальном отрезке,
удалённом от нулей и от частот.

\step{4} Пусть $\Pi_k=[\xi_k,\xi_{k+1}]\times[-\alpha,\alpha]$ и
\[
 P_k=-\frac1{2\pi i}\oint_{\partial\Pi_k}2\omega\,(T-\omega^2)^{-1}\,d\omega .
\]
Операторы $P_k$ идемпотентны, имеют конечный ранг, $P_jP_k=0$ при $j\ne k$, а их образы
$V_k$ лежат в $D(\Lambda)$ и инвариантны относительно $T$. Обозначим через $Q_k$
ортогональные спектральные проекторы оператора $\Lambda_c$ для тех же прямоугольников.
Из тождества для разности резольвент, оценки шага~3 и явных интегралов от
$\abs{\omega_i-\omega}^{-2}$ по сторонам прямоугольников получаем
\[
 \sum_k\abs{\inner{(P_k-Q_k)u}{v}}\le C\norm u\,\norm v ,
\]
откуда следуют двусторонние неравенства Рисса для $V_k$
(ср.~\cite[предложение~3.7]{Wyss}). Полнота системы $\{V_k\}$ получается из интеграла от
резольвенты по большому симметричному прямоугольнику.
\end{proof}

\begin{proof}[Доказательство теоремы~\ref{th:sign}]
Согласно леммам~\ref{lem:sub} и~\ref{lem:count} к паре $\Lambda$, $\tilde S$ применима
теорема~\ref{th:mm}. Положим $W_k=rV(V_k)$. Эти подпространства конечномерны, лежат в
$D(L_a)$, инвариантны относительно $L_a=rV(\Lambda+\tilde S)(rV)^{-1}$ и образуют базис
Рисса подпространств как образ такого базиса при ограниченно обратимом операторе.
Ограничение $L_a$ на $W_k$ есть линейный оператор в конечномерном комплексном
пространстве, поэтому $W_k$ --- сумма его корневых подпространств. Теорема доказана.
\end{proof}

\section{Метод подобных операторов}
Доказательство теоремы~\ref{th:mm} использует контурные интегралы и оценку минимума модуля
голоморфной функции. При малом $b_1$ подпространства $V_k$ строятся явно методом подобных
операторов (ср.~\cite{Pol}).

\begin{theorem}\label{th:small}
Пусть выполнены условия теоремы~\ref{th:mm} и $160\,b_1c_2<1$. Тогда найдутся число
$c>0$, последовательность $\xi_0<\xi_1<\dots$, $\xi_k\to\infty$, и ограниченно обратимый
оператор $U$, для которых подпространства $V_k=UH_k$ конечномерны, лежат в $D(\Lambda)$,
инвариантны относительно $\Lambda+S$ и образуют базис Рисса подпространств. Здесь $H_k$
--- линейная оболочка векторов $e_i$ с $\xi_k<\sqrt{\lambda_i+c}<\xi_{k+1}$. Более того,
$\norm{U-I}\le1/5$.
\end{theorem}
\begin{proof}
Доказательство разобьём на четыре шага. Сохраняем обозначения $\Lambda_c$, $\omega_i$,
$\hat S$, $\beta$ шага~1 предыдущего доказательства; тогда
$T=\Lambda_c+\hat S\Lambda_c^{1/2}$.

\step{1} \emph{Разрезы.} Положим $\delta=1/(4c_2)$. В окне $[\tau,\tau+2\delta N]$ с
$N>2c_2\tau+2C_0$ среди точек $\tau+\delta(2q+1)$, $0\le q<N$, есть точка на расстоянии
не меньше $\delta$ от всех частот: иначе ниже $\tau+2\delta N$ лежало бы не менее $N$
частот, что противоречит условию~(i). Выбираем разрезы $\xi_k$ с таким зазором и
$\xi_{k+1}\ge2\xi_k$, причём $\xi_0$ лежит ниже всех частот. Блоки $H_k$ конечномерны и
попарно ортогональны, их сумма плотна; размерности блоков не ограничены в совокупности.

\step{2} \emph{Мультипликаторы.} Для ограниченного оператора $X$ с матрицей $(x_{il})$ в
базисе $\{e_i\}$ обозначим через $JX$ его блочно-диагональную часть, а через
$\Gamma_0X$ и $\Gamma_1X$ --- операторы, матрицы которых вне диагональных блоков имеют
элементы
\[
 \frac{\omega_l}{\omega_i^2-\omega_l^2}\,x_{il},\qquad
 \frac{\omega_i}{\omega_i^2-\omega_l^2}\,x_{il},
\]
а внутри блоков равны нулю. Покажем, что $\norm{\Gamma_jX}\le(4/\delta)\norm X$. Оба
множителя --- полуразность и полусумма чисел $1/(\omega_i-\omega_l)$ и
$1/(\omega_i+\omega_l)$. Если $\omega_i\ge v$ для строк и $\omega_l\le u<v$ для столбцов,
то $1/(\omega_i\mp\omega_l)=\sum_n(\pm\omega_l)^n\omega_i^{-n-1}$ есть сумма произведений
диагональных операторов, и умножение на этот множитель ограничено числом $1/(v-u)$. Пары
блоков с одной разностью номеров имеют ортогональные строки и столбцы, поэтому норма их
суммы равна максимуму норм. Удвоение разрезов даёт для разности номеров $n+1$ оценку
$2^{-n}/\delta$, и суммирование по $n$ завершает доказательство оценки.

\step{3} \emph{Уравнение.} Ищем оператор $U=I-\Gamma_0X$, для которого
\[
 TU=U\bigl(\Lambda_c+JX\,\Lambda_c^{1/2}\bigr).
\]
Из определений получаем
$\Lambda_c\Gamma_0X-\Gamma_0X\Lambda_c=(X-JX)\Lambda_c^{1/2}$ и
$\Gamma_1X\,\Lambda_c^{1/2}=\Lambda_c^{1/2}\Gamma_0X$. Поэтому искомое равенство
равносильно уравнению
\begin{equation}\label{eq:fixed}
 X=\hat S-\hat S\,\Gamma_1X+\Gamma_0X\cdot JX .
\end{equation}
В силу шага~2 правая часть~\eqref{eq:fixed} переводит шар $\norm X\le2\beta$ в себя и
является сжатием с постоянной $20\beta/\delta$, если $40\beta\le\delta$. Неравенство
$40(b_1+b_0/\sqrt c)\le1/(4c_2)$ выполнено при достаточно большом $c$ в том и только
том случае, когда $160\,b_1c_2<1$. Фиксируем такое $c$ и решение $X$
уравнения~\eqref{eq:fixed}.

\step{4} \emph{Подобие.} Имеем $\norm{\Gamma_0X}\le(4/\delta)\cdot2\beta\le1/5$, поэтому
оператор $U$ обратим. Оператор $\Lambda_c+JX\,\Lambda_c^{1/2}$ сохраняет каждый блок
$H_k$, а потому подпространства $V_k=UH_k$ инвариантны относительно $T$. Они образуют
базис Рисса подпространств как образ ортогонального разложения при ограниченно обратимом
операторе. Теорема доказана.
\end{proof}

\begin{corollary}\label{cor:near}
Пусть выполнены условия теоремы~\ref{th:sign} и почти всюду
\[
 \Bigl\lvert\frac{a'(x)}{a(x)}+\frac{a'(-x)}{a(-x)}\Bigr\rvert\le\eta,\qquad
 160\,(1+\sqrt\kappa)\,\eta<\pi\,(1-\sqrt\kappa).
\]
Тогда заключение теоремы~\ref{th:sign} следует из теоремы~\ref{th:small}.
\end{corollary}
\begin{proof}
В лемме~\ref{lem:gauge} положим $r(x)=\bigl(\abs{a(x)}/\abs{a(-x)}\bigr)^{1/4}$.
Коэффициент $b=\sigma\sqrt{a(x)a(-x)}$ остаётся прежним, а
\[
 \frac{r'(x)}{r(x)}=\frac14\Bigl(\frac{a'(x)}{a(x)}+\frac{a'(-x)}{a(-x)}\Bigr),
\]
так что в~\eqref{eq:sub} имеем $b_1\le\nu(1+\sqrt\kappa)\eta/2$. С постоянной
$c_2=2\nu/\pi$ из леммы~\ref{lem:count} получаем
$160\,b_1c_2\le160(1+\sqrt\kappa)\eta/(\pi(1-\sqrt\kappa))<1$. Следствие доказано.
\end{proof}

\begin{remark}\label{rem:near}
Существо условия следствия~\ref{cor:near} состоит в следующем. Запишем
$a=\sigma se^{\varphi}$ с чётной функцией $s>0$ и нечётной $\varphi$; тогда левая часть
условия равна $2\abs{\varphi'(x)}$. Таким образом, мала производная нечётной части
функции $\ln\abs a$, а на чётную часть ограничений нет; для чётного $a$ левая часть
равна нулю. Условие выполнено, например, для $a(x)=\tfrac12+\tfrac x{1000}$. Отметим, что
разрезы в теореме~\ref{th:small} удваиваются, поэтому размерности блоков этим методом не
оцениваются.
\end{remark}

\section{Полнота и суммируемость спектральных разложений}
\begin{theorem}\label{th:complete}
Пусть выполнено условие~\eqref{eq:bound} и функция $a$ непрерывна на отрезке $[-1,1]$.
Тогда система корневых функций оператора $L_a$ полна в $H$. Более того, найдутся
постоянные $C$ и $s_0>0$, для которых при всех $s\ge s_0$ оператор $L_a+s$ взаимно
однозначно отображает $D(L_a)$ на $H$ и
\[
 \norm{(L_a+s)^{-1}}\le\frac Cs .
\]
\end{theorem}
Знак $a$ и гладкость выше непрерывности здесь не используются. Условиям теоремы
удовлетворяют, в частности, коэффициенты $a(x)=\max(x,0)$ и $a(x)=x/2$, не охваченные
теоремами~\ref{th:even} и~\ref{th:sign}.

Идея доказательства состоит в следующем. Обозначим через $a_+$ и $a_-$ чётную и нечётную
части функции $a$:
\[
 a_\pm(x)=\tfrac12\bigl(a(x)\pm a(-x)\bigr).
\]
Если $\norm{a_+}_\infty<1$, то $\operatorname{Re}\inner{T_aW}{W}\ge c\norm W^2$ с
постоянной $c>0$, и резольвента оператора $L_a$ убывает на отрицательной полуоси как
$1/s$ (лемма~\ref{lem:coercive}). Сразу применить это соображение в условиях
теоремы~\ref{th:complete} затруднительно, так как чётная часть может быть большой: для
$a(x)=4\max(x,0)$ она равна $2\abs x$. Однако умножение на вес $e^{\theta/2}$ с нечётной
функцией $\theta$ заменяет $a$ коэффициентом $ae^{-\theta}$ и добавляет возмущение первого
порядка, а функцию $\theta$ можно выбрать так, чтобы чётная часть нового коэффициента
была по модулю меньше единицы (лемма~\ref{lem:phase}). В приведённом примере годится
$\theta(x)=2x$: чётная часть функции $ae^{-\theta}$ равна $2\abs xe^{-2\abs x}\le1/e$.

\begin{lemma}\label{lem:coercive}
Пусть $a$ --- вещественная функция из $L_\infty(-1,1)$ и $\norm{a_+}_\infty\le q<1$.
Тогда оператор $T_a$ ограниченно обратим, $\norm{T_a^{-1}}\le(1-q)^{-1}$, и при каждом
$s>0$ оператор $L_a+s$ взаимно однозначно отображает $D(L_a)$ на $H$, причём
\begin{equation}\label{eq:ray}
 \norm{(L_a+s)^{-1}}\le\frac{\norm{T_a}^2}{(1-q)^2\,s}.
\end{equation}
Более того, каждое собственное значение $\lambda$ оператора $L_a$ удовлетворяет
неравенствам
\[
 \operatorname{Re}\lambda>0,\qquad
 (1-q)\,\abs{\operatorname{Im}\lambda}\le\norm{a_-}_\infty\operatorname{Re}\lambda .
\]
\end{lemma}
Условие~\eqref{eq:bound} в лемме не предполагается.
\begin{proof}
Имеем $(aR)^*=Ra$ и $\bigl((aR\pm Ra)W\bigr)(x)=2a_\pm(x)W(-x)$. Поэтому для $W\in H$
\begin{align*}
 \abs{\operatorname{Re}\inner{aRW}{W}}&=\tfrac12\abs{\inner{(aR+Ra)W}{W}}\le q\norm W^2,\\
 \abs{\operatorname{Im}\inner{aRW}{W}}&=\tfrac12\abs{\inner{(aR-Ra)W}{W}}
 \le\norm{a_-}_\infty\norm W^2 ,
\end{align*}
откуда
\begin{equation}\label{eq:coercive}
 \operatorname{Re}\inner{T_aW}{W}\ge(1-q)\norm W^2 .
\end{equation}
Из~\eqref{eq:coercive} следует $\norm{T_aW}\ge(1-q)\norm W$, так что оператор $T_a$
инъективен и его образ замкнут. Если вектор $v$ ортогонален образу, то
$\inner{T_av}{v}=0$ и $v=0$ в силу~\eqref{eq:coercive}. Значит, $T_a$ обратим и
$\norm{T_a^{-1}}\le(1-q)^{-1}$. Оператор $G$ положителен, поэтому при $s\ge0$ те же
рассуждения применимы к оператору $T_a+sG$.

Пусть $f\in D(L_a)$ и $W=A_0f$. Тогда $f=GW$ и $(L_a+s)f=(T_a+sG)W$. Следовательно, при
$s\ge0$ оператор $L_a+s$ взаимно однозначно отображает $D(L_a)$ на $H$, и обратный к нему
равен $G(T_a+sG)^{-1}$.

Пусть теперь $s>0$, $(L_a+s)f=h$ и $u=T_a^{-1}f$. Применив $T_a^{-1}$ к равенству
$T_aW+sf=h$, получим $W+su=T_a^{-1}h$. Умножим это равенство скалярно на $f$. Так как
$\inner Wf=\norm{f'}^2\ge0$, имеем
\[
 s\operatorname{Re}\inner uf\le\operatorname{Re}\inner{T_a^{-1}h}{f}
 \le\frac{\norm h\,\norm f}{1-q}.
\]
С другой стороны, в силу~\eqref{eq:coercive}
\[
 \operatorname{Re}\inner uf=\operatorname{Re}\inner{T_au}{u}\ge(1-q)\norm u^2
 \ge\frac{(1-q)\norm f^2}{\norm{T_a}^2},
\]
и из двух последних неравенств получаем~\eqref{eq:ray}.

Наконец, пусть $L_af=\lambda f$, $f\ne0$ и $W=A_0f$. Тогда $T_aW=\lambda f$ и
$\inner{T_aW}{W}=\lambda E$, где $E=\inner fW=\norm{f'}^2>0$. Отделяя вещественную и
мнимую части, получаем
\[
 E\operatorname{Re}\lambda\ge(1-q)\norm W^2>0,\qquad
 E\,\abs{\operatorname{Im}\lambda}\le\norm{a_-}_\infty\norm W^2 ,
\]
откуда следует последнее утверждение. Лемма доказана.
\end{proof}
Для нечётного $a$ лемма применима с $q=0$ при любой величине $\norm a_\infty$:
собственные значения лежат в секторе
$\abs{\operatorname{Im}\lambda}\le\norm a_\infty\operatorname{Re}\lambda$. Для чётного
$a$ с $\abs a\le q<1$ имеем $a_-=0$, и собственные значения положительны, что согласуется
с теоремой~\ref{th:even}.

\begin{lemma}\label{lem:phase}
Пусть $a$ --- вещественная непрерывная функция на отрезке $[-1,1]$ и
$\abs{a(x)a(-x)}<1$ при всех $x$. Тогда найдутся нечётная функция
$\theta\in C^\infty[-1,1]$, равная нулю в окрестности нуля, и число $q<1$, для которых
\begin{equation}\label{eq:phase}
 \bigl\lvert a(x)e^{-\theta(x)}+a(-x)e^{\theta(x)}\bigr\rvert\le2q,\qquad x\in[-1,1].
\end{equation}
\end{lemma}
\begin{proof}
Доказательство разобьём на три шага. При $x\in[0,1]$ и $t\in\R$ положим
\[
 F(x,t)=a(x)e^{-t}+a(-x)e^{t},\qquad I(x)=\{t\in\R:\ \abs{F(x,t)}<2\}.
\]

\step{1} Покажем, что $I(x)$ --- непустой интервал. Пусть $p=a(x)$ и $\tilde p=a(-x)$;
по условию $\abs{p\tilde p}<1$. Если $p\tilde p>0$, то
$\abs{F(x,t)}=\abs pe^{-t}+\abs{\tilde p}e^{t}$ --- выпуклая функция переменной $t$ с
наименьшим значением $2\sqrt{p\tilde p}<2$, а множество, на котором выпуклая функция
меньше $2$, есть интервал. Если $p\tilde p<0$, то производная $-pe^{-t}+\tilde pe^{t}$
при всех $t$ имеет знак числа $\tilde p$, функция $F(x,\cdot)$ строго монотонна и
принимает все вещественные значения, поэтому прообраз интервала $(-2,2)$ --- непустой
интервал. Если $p\tilde p=0$, то функция $F(x,\cdot)$ равна $pe^{-t}$ или
$\tilde pe^{t}$, и $I(x)$ --- полупрямая или вся прямая.

\step{2} Так как $\abs{F(0,0)}=2\abs{a(0)}<2$, найдётся $\delta\in(0,\tfrac12)$, для
которого $\abs{F(x,0)}<2$ при $x\in[0,2\delta]$. Положим $U_0=(-1,2\delta)$ и $t_0=0$.
Для каждого $y\in[\delta,1]$ выберем $t_y\in I(y)$ и открытый интервал
$U_y\subset(\delta/2,+\infty)$, содержащий $y$, для которого $\abs{F(x,t_y)}<2$ при
$x\in U_y\cap[0,1]$; он существует в силу непрерывности функции $F(\cdot,t_y)$. Из
покрытия отрезка $[0,1]$ множествами $U_0$ и $U_y$ выберем конечное подпокрытие
$U_0,U_1,\dots,U_n$ с числами $t_0,t_1,\dots,t_n$. Множество $U_0$ в него входит, так как
точка $0$ лежит только в $U_0$. Если $x\in U_i\cap[0,1]$, то $t_i\in I(x)$.

\step{3} Пусть $\{\chi_i\}_{i=0}^n$ --- гладкое разбиение единицы на отрезке $[0,1]$,
подчинённое этому покрытию: $\chi_i\ge0$, $\operatorname{supp}\chi_i\subset U_i$ и
$\sum_i\chi_i=1$ на $[0,1]$. Положим $\theta(x)=\sum_i\chi_i(x)\,t_i$ при $x\in[0,1]$.
Число $\theta(x)$ есть выпуклая комбинация тех $t_i$, для которых $x\in U_i$. Все они
лежат в интервале $I(x)$, поэтому $\theta(x)\in I(x)$. При $x\in[0,\delta/2]$ точка $x$
лежит только в $U_0$, и $\theta(x)=0$. Продолжим $\theta$ на $[-1,0)$ равенством
$\theta(x)=-\theta(-x)$; полученная функция нечётна и бесконечно дифференцируема. Функция
$x\mapsto\abs{F(x,\theta(x))}$ непрерывна на $[0,1]$ и в каждой точке меньше $2$;
обозначим её максимум через $2q$. Тогда~\eqref{eq:phase} выполнено при $x\in[0,1]$, а при
замене $x$ на $-x$ левая часть~\eqref{eq:phase} не меняется. Лемма доказана.
\end{proof}

Напомним нужные понятия из~\cite[\S\S2,~4]{Shk}. Пусть $T$ --- плотно определённый оператор с
компактным обратным. Говорят, что $T$ имеет порядок $2$, если $j$-е сингулярное число
оператора $T^{-1}$ не превосходит $cj^{-2}$. Луч $\arg\lambda=\phi$ называется лучом
наилучшего убывания резольвенты, если на нём
$\norm{(T-\lambda)^{-1}}\le C\abs\lambda^{-1}$ при $\abs\lambda>r_0$.
Согласно~\cite[теорема~4.7]{Shk} система корневых векторов оператора порядка $2$,
имеющего хотя бы один такой луч, полна (в обозначениях~\cite{Shk} такой оператор
принадлежит классу $\mathcal K_2$).

\begin{proof}[Доказательство теоремы~\ref{th:complete}]
Доказательство разобьём на четыре шага.

\step{1} \emph{Калибровка.} Функция $a$ непрерывна, поэтому
неравенство~\eqref{eq:bound} выполнено в каждой точке отрезка $[-1,1]$. Пусть $\theta$ и
$q$ --- функция и число из леммы~\ref{lem:phase}. Положим $r=e^{\theta/2}$ и
$b=ae^{-\theta}$. Функции $r$ и $1/r$ принадлежат $W^2_\infty(-1,1)$, а из нечётности
$\theta$ следует $r(-x)=e^{-\theta(x)/2}$ и $b(x)r(x)=a(x)r(-x)$. По
лемме~\ref{lem:gauge} для $g\in D(A_0)$
\begin{equation}\label{eq:phasegauge}
 L_a(rg)=r\bigl(L_bg+T_bB_0g\bigr),\qquad
 B_0g=-\theta'g'-\varrho g,\qquad\varrho=\tfrac12\theta''+\tfrac14\theta'^2 .
\end{equation}
Функция $b$ вещественна и ограничена, а $2b_+(x)$ совпадает с выражением под знаком
модуля в~\eqref{eq:phase}. Следовательно, $\norm{b_+}_\infty\le q<1$.

\step{2} \emph{Подчинённость.} По лемме~\ref{lem:coercive} имеем
$\norm{T_b^{-1}}\le(1-q)^{-1}$, и при $s>0$ оператор $R_s=(L_b+s)^{-1}$ определён на всём
$H$ и $\norm{R_s}\le C_1/s$, где $C_1=\norm{T_b}^2(1-q)^{-2}$. Пусть $g\in D(A_0)$. Тогда
\[
 \norm{g'}^2=\inner{A_0g}{g}\le\norm{A_0g}\,\norm g,\qquad
 \norm g=\norm{GA_0g}\le\norm G\,\norm{A_0g},
\]
а из равенства $L_bg=T_bA_0g$ следует $\norm{A_0g}\le(1-q)^{-1}\norm{L_bg}$. Поэтому
\[
 \norm{T_bB_0g}\le\norm{T_b}\bigl(\norm{\theta'}_\infty\norm{g'}
 +\norm\varrho_\infty\norm g\bigr)\le C_2\,\norm g^{1/2}\norm{L_bg}^{1/2},
\]
где $C_2=\norm{T_b}\bigl(\norm{\theta'}_\infty+\norm\varrho_\infty\norm G^{1/2}\bigr)
(1-q)^{-1/2}$.

\step{3} \emph{Оценка резольвенты.} Пусть $s>0$, $h\in H$ и $g=R_sh$. Тогда
$\norm g\le C_1\norm h/s$ и $\norm{L_bg}=\norm{h-sg}\le(1+C_1)\norm h$. В силу шага~2
оператор $K_s=T_bB_0R_s$ ограничен и
\[
 \norm{K_s}\le C_2\sqrt{C_1(1+C_1)/s}.
\]
Выберем $s_0>0$ так, чтобы $\norm{K_s}\le\tfrac12$ при $s\ge s_0$. Тогда оператор
$I+K_s$ обратим и $\norm{(I+K_s)^{-1}}\le2$. На $D(A_0)$ имеем
$L_b+T_bB_0+s=(I+K_s)(L_b+s)$. Значит, оператор $L_b+T_bB_0+s$ взаимно однозначно
отображает $D(A_0)$ на $H$, и обратный к нему равен $R_s(I+K_s)^{-1}$. Умножение на $r$
взаимно однозначно отображает на себя и $D(A_0)$, и $H$. Отсюда и
из~\eqref{eq:phasegauge} следует, что оператор $L_a+s$ взаимно однозначно отображает
$D(L_a)$ на $H$ и
\[
 \norm{(L_a+s)^{-1}}=\bigl\lVert rR_s(I+K_s)^{-1}r^{-1}\bigr\rVert
 \le\frac{2C_1e^{\norm\theta_\infty}}{s},\qquad s\ge s_0 .
\]

\step{4} \emph{Полнота.} Оператор $L_a$ плотно определён, а обратный к нему
$K_a=GT_a^{-1}$ компактен по теореме~\ref{th:class}. Собственные значения оператора $G$
равны $(2/(j\pi))^2$, поэтому $j$-е сингулярное число оператора $K_a$ не превосходит
$\norm{T_a^{-1}}(2/(j\pi))^2$, и $L_a$ имеет порядок $2$. В силу шага~3 отрицательная
полуось есть луч наилучшего убывания резольвенты оператора $L_a$. Остаётся применить
\cite[теорема~4.7]{Shk}. Теорема доказана.
\end{proof}

Если $\norm{a_+}_\infty<1$, то спектральные разложения по корневым функциям суммируются
методом Абеля--Лидского. Напомним определение из~\cite[\S5]{Shk}. Положим
\[
 \Lambda_\vartheta=\{\lambda\in\C:\ \abs{\arg\lambda}\le\vartheta\pi/2\},\qquad
 0\le\vartheta<2 .
\]
Пусть $T$ --- плотно определённый оператор с компактным обратным и $\beta>0$. Будем
говорить, что система корневых векторов оператора $T$ образует \emph{базис для метода
суммирования Абеля--Лидского порядка~$\beta$}, если найдутся положительно ориентированные
замкнутые контуры $\Gamma_k$, каждый из которых охватывает конечное число собственных
значений $T$, такие что при каждом $t>0$ и каждом $x\in H$ ряд
\[
 S(t)x=-\sum_k\frac1{2\pi i}\int_{\Gamma_k}e^{-\lambda^\beta t}(T-\lambda)^{-1}x\,d\lambda
\]
сходится и $S(t)x\to x$ при $t\to0+$. Здесь $\lambda^\beta$ --- главная ветвь степени.
Контуры $\Gamma_k$ от $x$ и $t$ не зависят; их выбор равносилен расстановке скобок в ряде
из проекторов Рисса на корневые подпространства.

\begin{lemma}\label{lem:sector}
Пусть $a$ --- вещественная функция из $L_\infty(-1,1)$ и $\norm{a_+}_\infty\le q<1$.
Положим
\[
 \tau=\frac2\pi\arctan\frac{\norm{a_-}_\infty}{1-q}.
\]
Пусть $0<\delta\le\pi/2$ и $\abs{\arg\lambda}\ge\pi\tau/2+\delta$. Тогда $\lambda$
принадлежит резольвентному множеству оператора $L_a$ и
\begin{equation}\label{eq:sector}
 \abs\lambda\,\norm{(L_a-\lambda)^{-1}}\le\frac{\norm{T_a}}{(1-q)\sin\delta}+1 .
\end{equation}
\end{lemma}
\begin{proof}
Пусть $f\in D(L_a)$, $(L_a-\lambda)f=g$ и $W=A_0f$. Тогда $T_aW-\lambda f=g$. Умножим это
равенство скалярно на $W$. Так как $\inner fW=\norm{f'}^2=E\ge0$, получаем
\[
 z-\lambda E=\inner gW,\qquad z=\inner{T_aW}{W}.
\]
В доказательстве леммы~\ref{lem:coercive} установлено, что
$\operatorname{Re}z\ge(1-q)\norm W^2$ и
$\abs{\operatorname{Im}z}\le\norm{a_-}_\infty\norm W^2$. Поэтому
$\abs{\arg z}\le\pi\tau/2$ и $\abs z\ge(1-q)\norm W^2$. Точка $\lambda E$ лежит на луче
$\arg\zeta=\arg\lambda$, а угол между этим лучом и направлением на $z$ не меньше $\delta$.
Расстояние от $z$ до луча равно $\abs z\sin\varphi$, если угол $\varphi$ между ними не
больше $\pi/2$, и равно $\abs z$ в противном случае. Следовательно,
\[
 (1-q)\sin\delta\,\norm W^2\le\abs{z-\lambda E}\le\norm g\,\norm W,
\]
то есть $\norm W\le\norm g/((1-q)\sin\delta)$. Из равенства $\lambda f=T_aW-g$ получаем
$\abs\lambda\norm f\le\norm{T_a}\norm W+\norm g$, откуда следует~\eqref{eq:sector} для
каждого решения. При $g=0$ отсюда $f=0$, так как $\lambda\ne0$. Обратный к $L_a$ оператор
$GT_a^{-1}$ компактен по лемме~\ref{lem:coercive}, поэтому спектр $L_a$ состоит из
собственных значений, и $\lambda$ лежит в резольвентном множестве. Лемма доказана.
\end{proof}

\begin{theorem}\label{th:abel}
Пусть $a$ --- вещественная функция из $L_\infty(-1,1)$, $\norm{a_+}_\infty\le q<1$, и
число $\tau$ определено в лемме~\ref{lem:sector}. Тогда система корневых функций
оператора $L_a$ образует базис для метода суммирования Абеля--Лидского порядка $\beta$
при любом $\beta\in(1/2,\tau^{-1})$, а при $\tau=0$ --- при любом $\beta>1/2$.
\end{theorem}
\begin{proof}
Проверим условия теоремы Лидского--Мацаева \cite[теорема~5.1]{Shk}. Оператор $L_a$ плотно
определён и имеет компактный обратный $GT_a^{-1}$, а $j$-е сингулярное число этого
обратного не превосходит $\norm{T_a^{-1}}(2/(j\pi))^2$; значит, $L_a$ имеет порядок~$2$.
По лемме~\ref{lem:coercive} все собственные значения, а с ними весь спектр, лежат в
секторе $\Lambda_\tau$, и $\tau<1<2$. Каждый луч вне $\Lambda_\tau$ имеет вид
$\arg\lambda=\phi$ с $\abs\phi>\pi\tau/2$, и по лемме~\ref{lem:sector} он является лучом
наилучшего убывания резольвенты. По теореме~5.1 из~\cite{Shk} корневые векторы образуют
базис для метода Абеля--Лидского порядка $\beta$ при $\beta\in(1/2,\tau^{-1})$.
Отметим, что доказательство этой теоремы использует оценку
$\norm{(L_a-\lambda)^{-1}}\le C\abs\lambda^{-1}$ с одной постоянной $C$ во всём
дополнении сектора $\Lambda_{\tau+\varepsilon}$, $\varepsilon>0$; оценка~\eqref{eq:sector}
с $\delta=\min(\pi\varepsilon/2,\pi/2)$ её даёт. Теорема доказана.
\end{proof}
Для $a(x)=x/2$ имеем $q=0$ и $\tau=\frac2\pi\arctan\frac12$, так что допустим любой
порядок $\beta\in(1/2,\tau^{-1})$, $\tau^{-1}\approx3{,}39$. Для $a(x)=\max(x,0)$ имеем
$a_+(x)=\abs x/2$, $a_-(x)=x/2$, $q=1/2$, $\tau=1/2$ и $\beta\in(1/2,2)$. В обоих случаях
базиса Рисса из собственных функций нет (теоремы~\ref{th:ramp} и~\ref{th:half}).

\begin{remark}\label{rem:complete}
Непрерывность в теореме~\ref{th:complete} отбросить нельзя: разрывный коэффициент
теоремы~\ref{th:onesided} удовлетворяет условию~\eqref{eq:bound}, а собственных функций у
оператора $L_a$ нет. Из полноты базисность не следует: в теоремах~\ref{th:ramp}
и~\ref{th:half} коэффициент непрерывен, а базиса Рисса из собственных функций нет.
Отметим, что в условиях теоремы~\ref{th:complete} спектр исключён только на луче
$(-\infty,-s_0]$; сектор, содержащий все собственные значения, и суммируемость по
Абелю--Лидскому получены лишь при $\norm{a_+}_\infty<1$ (лемма~\ref{lem:coercive},
теорема~\ref{th:abel}).
\end{remark}

\section{Примеры, в которых базисности нет}\label{sec:examples}
Покажем, что условие~\eqref{eq:bound} само по себе базисности не даёт. В первых двух
примерах $a(x)a(-x)\equiv0$, так что~\eqref{eq:bound} выполнено с $\kappa=0$, а
$(aR)^2=0$: половины интервала связаны односторонне. В третьем примере коэффициент
нечётен и аналитичен.

\begin{theorem}\label{th:onesided}
Пусть $a(x)=4$ при $0<x<1$ и $a(x)=0$ при $-1<x\le0$. Тогда при любом $\lambda\in\C$
уравнение $L_ay=\lambda y$ имеет в $D(L_a)$ только нулевое решение. Более того, не
существует базиса Рисса подпространств, состоящего из конечномерных подпространств,
инвариантных относительно $L_a^{-1}$.
\end{theorem}
\begin{proof}
Пусть $u(x)=y(x)$ и $v(x)=y(-x)$ на $(0,1)$. Тогда
\[
 v''=-\lambda v,\qquad u''=-\lambda u+4\lambda v,\qquad
 u(0)=v(0),\quad u'(0)=-v'(0),\quad u(1)=v(1)=0 .
\]
Функция $d=u+2xv'-v$ удовлетворяет уравнению $d''=-\lambda d$ и условиям
$d(0)=d'(0)=0$, поэтому $d\equiv0$, то есть $u=v-2xv'$. При $x=1$ получаем $v'(1)=0$;
вместе с $v(1)=0$ это даёт $v\equiv0$ и $u\equiv0$. Наконец, ненулевое конечномерное
инвариантное подпространство оператора $L_a^{-1}$ содержало бы собственный вектор.
Теорема доказана.
\end{proof}
Коэффициент теоремы~\ref{th:onesided} разрывен. В следующем примере он липшицев.

\begin{theorem}\label{th:ramp}
Пусть $a(x)=\max(x,0)$. Тогда не существует базиса Рисса пространства $H$, составленного
из собственных функций оператора $L_a$.
\end{theorem}
\begin{proof}
Доказательство разобьём на три шага.

\step{1} Пусть $L_ay=\lambda y$, $y\ne0$, $\abs\lambda\ge\tfrac18$ и
$\bigl\lvert\abs\lambda-\lambda\bigr\rvert\le\gamma\sqrt{\abs\lambda}$. Покажем, что
$\abs\lambda\le32e^{2\gamma}$. Для решения уравнения $w''=-\lambda w$ положим
$E_w=\abs{w'}^2+\abs\lambda\abs w^2$; тогда $\abs{E_w'}\le\gamma E_w$, и на отрезке
длины $1$ величина $E_w$ растёт не более чем в $e^\gamma$ раз. На $[-1,0]$ функция $y$
сама удовлетворяет уравнению $y''=-\lambda y$, откуда
$E_y(0)\le e^\gamma\abs{y'(-1)}^2$. На $(0,1)$ для $u(x)=y(x)$, $v(x)=y(-x)$ имеем
$v''=-\lambda v$, $u''=-\lambda u+\lambda xv$, и функция
\[
 p=u+\frac{x^2}4\,v'-\frac x4\,v
\]
удовлетворяет уравнению $p''=-\lambda p$ и равенствам $p(0)=y(0)$,
$p'(0)=y'(0)-y(0)/4$, $p(1)=-y'(-1)/4$. Поэтому
\[
 \frac{\abs\lambda}{16}\abs{y'(-1)}^2\le E_p(1)\le e^\gamma E_p(0)\le2e^\gamma E_y(0)
 \le2e^{2\gamma}\abs{y'(-1)}^2 .
\]
Если $y'(-1)=0$, то $y\equiv0$ в силу единственности решения задачи Коши.

\step{2} Пусть $\psi$ --- гладкая функция, $\psi=0$ при $x\le0$, $\psi(1)=0$,
$\psi\not\equiv0$, и $y_\omega=\psi(x)e^{i\omega x}$, $\omega>0$. Тогда
$y_\omega\in D(L_a)$, $a(x)y_\omega''(-x)\equiv0$ и
\[
 \norm{(L_a-\omega^2)y_\omega}=\norm{2i\omega\psi'+\psi''}\le A+B\omega,\qquad
 \norm{y_\omega}=\norm\psi .
\]

\step{3} Пусть $\{e_i\}$ --- базис Рисса из собственных функций и
$L_ae_i=\lambda_ie_i$. Если $\abs{\lambda_i-\mu}\ge\delta$ при всех $i$, то из неравенств
Рисса следует $\norm{(L_a-\mu)y}\ge\delta\norm y/C$ для $y\in D(L_a)$. Согласно шагу~2
при каждом $\omega$ найдётся $\lambda_i$ с $\abs{\lambda_i-\omega^2}\le D_0+D_1\omega$.
Для такого $\lambda_i$ при больших $\omega$ выполнены условия шага~1 с
$\gamma=4(D_0+D_1)$, но $\abs{\lambda_i}\to\infty$. Полученное противоречие завершает
доказательство. Теорема доказана.
\end{proof}
В теоремах~\ref{th:onesided} и~\ref{th:ramp} коэффициент равен нулю на половине
интервала. В следующем примере он аналитичен и нечётен.

\begin{theorem}\label{th:half}
Пусть $a(x)=x/2$. Тогда не существует базиса Рисса пространства $H$, составленного из
собственных функций оператора $L_a$.
\end{theorem}
Здесь условие~\eqref{eq:bound} выполнено с $\kappa=\tfrac14$. По
теореме~\ref{th:complete} система корневых функций полна, а по
лемме~\ref{lem:coercive} собственные значения лежат в секторе
$\abs{\operatorname{Im}\lambda}\le\tfrac12\operatorname{Re}\lambda$.

Идея доказательства состоит в следующем. Для нечётного $a$ оператор $L_a$ унитарно
эквивалентен оператору $(1+ia)A_0$. При $a(x)=x/2$ находим функции $y_t$, $t\ge0$,
сосредоточенные около нуля, для которых $\norm{(L_a-t^2)y_t}\le C_0\norm{y_t}$ с
постоянной $C_0$, не зависящей от $t$. Если собственные функции образуют базис Рисса, то
около каждой точки $t^2$ на ограниченном расстоянии есть собственное значение, и число
собственных значений в круге радиуса $\ell$ растёт не медленнее $\ell$. Но сингулярные
числа оператора $L_a^{-1}$ допускают лишь рост порядка $\sqrt\ell$.

\begin{proof}
Доказательство разобьём на пять шагов.

\step{1} Пусть $a$ --- нечётная функция из $L_\infty(-1,1)$ и $U=P_++iP_-$. Оператор $U$
унитарен, сохраняет $D(A_0)$ и коммутирует с $A_0$. Умножение на нечётную функцию меняет
чётность, а $T_aW=W+a(P_+W-P_-W)$, поэтому $P_+T_aW=P_+W-aP_-W$,
$P_-T_aW=P_-W+aP_+W$ и
\[
 UT_aW=(1+ia)P_+W+i(1+ia)P_-W=(1+ia)\,UW,\qquad W\in H .
\]
Следовательно, $UL_a=(1+ia)A_0U$ на $D(L_a)$.

\step{2} Пусть $a(x)=x/2$ и $\zeta(x)=1+ix/2$. При $t\ge0$ положим
$u_t=\zeta^{1/4}\exp(-4t\zeta^{1/2})$ с главными значениями степеней. Дифференцируя,
получаем
\[
 -\zeta u_t''-t^2u_t=-\frac3{64\,\zeta}\,u_t,\qquad
 u_t'=\frac i2\Bigl(\frac1{4\zeta}-\frac{2t}{\zeta^{1/2}}\Bigr)u_t .
\]
Далее, $\abs{u_t}=\abs\zeta^{1/4}\exp(-4t\operatorname{Re}\zeta^{1/2})$ и
$(\operatorname{Re}\zeta^{1/2})^2=(\abs\zeta+1)/2$, где $\abs\zeta^2=1+x^2/4$. Отсюда
\[
 \operatorname{Re}\zeta^{1/2}\le1+\tfrac1{512}\quad\text{при }\abs x\le\tfrac14,\qquad
 \operatorname{Re}\zeta^{1/2}\ge1+\tfrac1{200}\quad\text{при }\abs x\ge\tfrac12 .
\]

\step{3} Пусть $\chi$ --- гладкая функция на $[-1,1]$, $\chi=1$ при
$\abs x\le\tfrac12$, $\chi(\pm1)=0$, и $Z_t=\chi u_t$. Тогда $Z_t\in D(A_0)$ и
\[
 (\zeta A_0-t^2)Z_t=-\frac3{64\,\zeta}\,Z_t-\zeta\,(2\chi'u_t'+\chi''u_t).
\]
Последнее слагаемое равно нулю при $\abs x<\tfrac12$. Так как $1\le\abs\zeta<2$, из
шага~2 получаем $\abs{u_t'}\le(1+t)\abs{u_t}$. Поэтому норма этого слагаемого не
превосходит $C'(1+t)e^{-4t(1+1/200)}$ с постоянной $C'$, не зависящей от $t$. С другой
стороны,
\[
 \norm{Z_t}^2\ge\int_{-1/4}^{1/4}\abs{u_t}^2dx\ge\tfrac12\,e^{-8t(1+1/512)} .
\]
Поэтому $\norm{(\zeta A_0-t^2)Z_t}\le C_0\norm{Z_t}$ с постоянной $C_0$, не зависящей от
$t\ge0$. Положим $y_t=U^{-1}Z_t$. Согласно шагу~1
$\norm{(L_a-t^2)y_t}=\norm{(\zeta A_0-t^2)Z_t}\le C_0\norm{y_t}$.

\step{4} Пусть $\{e_i\}$ --- базис Рисса из собственных функций, $L_ae_i=\lambda_ie_i$ и
\[
 m_0\sum_i\abs{c_i}^2\le\Bigl\lVert\sum_ic_ie_i\Bigr\rVert^2\le M_0\sum_i\abs{c_i}^2
\]
при любом конечном наборе чисел $c_i$. Повторяя рассуждение шага~3 доказательства
теоремы~\ref{th:ramp} с функциями $y_t$, получаем, что при каждом $t\ge0$ найдётся
$\lambda_i$ с $\abs{\lambda_i-t^2}\le D$, где $D$ не зависит от $t$.
Круги радиуса $D$ с центрами $3Dk$, $k=1,\dots,n$, попарно не пересекаются. Значит, при
каждом $n$ найдётся $n$ различных индексов $i$, для которых $\abs{\lambda_i}\le4Dn$.

\step{5} Покажем, что число индексов $i$ с $\abs{\lambda_i}\le\ell$ не превосходит
$c_0\sqrt\ell+2$ с постоянной $c_0$, не зависящей от $\ell$. Имеем $\lambda_i\ne0$ и
$K_ae_i=\lambda_i^{-1}e_i$. Как в доказательстве леммы~\ref{lem:count}, запишем
$G=F_N+E_N$, где $\operatorname{rank}F_N\le2N$ и $\norm{E_N}<(N\pi)^{-2}$, и выберем
наименьшее натуральное $N$, для которого
$\norm{T_a^{-1}}(N\pi)^{-2}\le(m_0/M_0)^{1/2}\ell^{-1}$. Пусть нашлось более $2N$ индексов
с $\abs{\lambda_i}\le\ell$. Тогда в линейной оболочке соответствующих $e_i$ есть
ненулевой вектор $u=\sum_ic_ie_i$ с $F_NT_a^{-1}u=0$, и
\[
 \Bigl(\frac{m_0}{M_0}\Bigr)^{1/2}\frac{\norm u}{\ell}\le
 \Bigl(m_0\sum_i\abs{c_i}^2\abs{\lambda_i}^{-2}\Bigr)^{1/2}\le\norm{K_au}
 =\norm{E_NT_a^{-1}u}<\frac{\norm{T_a^{-1}}}{(N\pi)^2}\,\norm u ,
\]
что противоречит выбору $N$. Значит, таких индексов не больше $2N\le c_0\sqrt\ell+2$.
Согласно шагу~4 получаем $n\le c_0\sqrt{4Dn}+2$ при всех $n$, что неверно при больших $n$.
Теорема доказана.
\end{proof}

\begin{remark}\label{rem:brackets}
Пусть $\{W_k\}$ --- базис Рисса подпространств и в каждом $W_k$ выбран базис
$\{e_{k,j}\}_j$. Объединённая система образует базис Рисса в $H$ в том и только том
случае, когда найдутся общие для всех $k$ постоянные $0<m_0\le M_0$, для которых
\[
 m_0\sum_j\abs{c_j}^2\le\Bigl\lVert\sum_jc_je_{k,j}\Bigr\rVert^2\le M_0\sum_j\abs{c_j}^2
\]
(ср.~\cite[предложение~3.2]{Wyss}). Для подпространств теоремы~\ref{th:sign} такая
равномерность не установлена: ни их размерности, ни обусловленность корневых функций
внутри них не оценены. Нам неизвестно, образуют ли корневые функции оператора $L_a$
обычный базис Рисса в условиях теоремы~\ref{th:sign} или хотя бы
следствия~\ref{cor:near}.
\end{remark}

\begin{remark}\label{rem:open}
Теоремы~\ref{th:ramp} и~\ref{th:half} исключают базис Рисса из собственных функций,
тогда как по теоремам~\ref{th:complete} и~\ref{th:abel} система корневых функций в
обоих примерах полна и образует базис для метода Абеля--Лидского порядка $\beta$ при
$\beta\in(1/2,\tau^{-1})$. Порядок $\beta\le1/2$ мы не исследовали. Нам неизвестно, существует ли для коэффициентов $a(x)=\max(x,0)$ и $a(x)=x/2$ базис
Рисса со скобками или базис, включающий присоединённые функции. Метод \S4 к
коэффициенту, меняющему знак, неприменим: в лемме~\ref{lem:gauge}
$b(x)b(-x)=a(x)a(-x)$, и при $a(x)a(-x)<0$ вещественный коэффициент $b$ не может быть
чётным.
\end{remark}

\begin{thebibliography}{9}
\bibitem{Wyss}
C.~Wyss. \emph{Riesz bases for $p$-subordinate perturbations of normal operators}.
arXiv:0906.2631v1.
\bibitem{Shk}
А.~А.~Шкаликов. \emph{Возмущения самосопряженных и нормальных операторов с дискретным
спектром}. Успехи матем. наук 71:5(431) (2016), 113--174.
\bibitem{Pol}
D.~M.~Polyakov. \emph{On the Bari basis property for even-order differential operators
with involution}. Tamkang J. Math. 54(4) (2023), 339--351.
\bibitem{KS}
A.~Kopzhassarova, A.~Sarsenbi. \emph{Basis properties of eigenfunctions of second-order
differential operators with involution}. Abstr. Appl. Anal. 2012, Article 576843.
\end{thebibliography}
\end{document}
